\section{The tools}
\label{sec:tools}

The models of this thesis are built from three ingredients: a count that jumps by one,
a clock that decides when it jumps, and a generating function that turns an infinite
system of \de s into a single equation one can hope to solve. This section fixes each
in turn, in the order in which one needs them, and closes with the linear birth--death
process --- the continuous-time reference model against which every later construction
is measured.

\subsection{Exponential clocks and their competition}
\label{sec:exp}

Waiting times enter first, because everything else is assembled from them. A
non-negative \rv{} $T$ is exponentially distributed with rate $\lambda>0$, written
$T\sim\text{Exp}(\lambda)$, when it has density and distribution function
\begin{equation}
\label{eq:expdef}
f_T(x) = \lambda\ee^{-\lambda x},
\qquad
F_T(x) = \pr{T\le x} = 1-\ee^{-\lambda x}
\qquad (x\ge0),
\end{equation}
with $\mean{T} = 1/\lambda$ and $\text{Var}(T) = 1/\lambda^2$. The rate is the reciprocal
of the mean waiting time, and it is the parameter in which models are specified.

The property that earns the exponential law its place is that waiting confers no
advantage.

\begin{definition}[Memorylessness]
\label{def:memoryless}
A non-negative \rv{} $T$ is memoryless if
\begin{equation}
\label{eq:memoryless}
\pr{T > t_1+t_2 \mid T > t_1} = \pr{T > t_2}
\qquad\text{for all } t_1,t_2>0 .
\end{equation}
\end{definition}

\noindent
For $T\sim\text{Exp}(\lambda)$ this is immediate from \cref{eq:expdef}, since
$\pr{T>x} = \ee^{-\lambda x}$ and therefore
\begin{equation*}
\pr{T > t_1+t_2 \mid T > t_1}
= \frac{\ee^{-\lambda(t_1+t_2)}}{\ee^{-\lambda t_1}}
= \ee^{-\lambda t_2} = \pr{T > t_2}.
\end{equation*}
Among continuous laws the exponential is the only one with this property, which is why
a model whose event rates do not depend on elapsed time is an exponential-clock model
whether or not it was described as one.

Populations offer several events at once --- a birth, a death, an absorption, a
catastrophe --- so what matters in practice is not a single clock but a race between
them. Take independent clocks $T_1\sim\text{Exp}(\lambda_1)$ and
$T_2\sim\text{Exp}(\lambda_2)$. Conditioning on the value of $T_2$ and integrating,
\begin{equation}
\label{eq:race2}
\pr{T_1<T_2}
= \int_0^{\infty}\!\!\lrb{\int_0^{y}\lambda_1\ee^{-\lambda_1 x}\dd x}\lambda_2\ee^{-\lambda_2 y}\dd y
= \int_0^{\infty}\!\!\lrb{1-\ee^{-\lambda_1 y}}\lambda_2\ee^{-\lambda_2 y}\dd y
= \frac{\lambda_1}{\lambda_1+\lambda_2},
\end{equation}
while for the time of the first event
\begin{equation}
\label{eq:min2}
\pr{\min\{T_1,T_2\}>x} = \pr{T_1>x}\pr{T_2>x} = \ee^{-(\lambda_1+\lambda_2)x},
\end{equation}
so that $\min\{T_1,T_2\}\sim\text{Exp}(\lambda_1+\lambda_2)$. The same two calculations
run unchanged for $n$ independent clocks with rates $\lambda_1,\dots,\lambda_n$. Writing
$\Lambda = \sum_i\lambda_i$ for the total rate,
\begin{equation}
\label{eq:racen}
\min_i T_i \sim \text{Exp}(\Lambda),
\qquad
\pr{T_i = \min_j T_j} = \frac{\lambda_i}{\Lambda},
\end{equation}
and the two are independent: knowing which channel fired first tells one nothing about
when it fired. Rates therefore do two separate jobs, setting the pace of events through
$\Lambda$ and apportioning them among channels through the ratios $\lambda_i/\Lambda$.

\begin{remark}[Simulation]
\label{rem:gillespie}
\Cref{eq:racen} is also an algorithm. From the current state, compute the rate $a_k$ of
each available event and their total $\Lambda$; draw a holding time
$\tau\sim\text{Exp}(\Lambda)$; select event $k$ with probability $a_k/\Lambda$; update
the state and advance the clock by $\tau$. This is Gillespie's stochastic simulation
algorithm, and it is exact for the chains defined below. Simulation is used in this
chapter only to check analysis, and the figures say where.
\end{remark}

\subsection{Continuous-time Markov chains}
\label{sec:ctmc}

A stochastic process is a family $\{\xt : t\ge0\}$ of random variables on a common
probability space. Throughout this thesis the index set is $[0,\infty)$ and the state
space is a subset of $\{0,1,2,\dots\}$: time runs continuously, and the state is a
count. The process is a \emph{continuous-time Markov chain} (CTMC) when its future
depends on its past only through its present,
\begin{equation}
\label{eq:markov}
\pr{\xx_{t_{n+1}} = i_{n+1} \mid \xx_{t_0}=i_0,\dots,\xx_{t_n}=i_n}
= \pr{\xx_{t_{n+1}} = i_{n+1} \mid \xx_{t_n}=i_n}
\end{equation}
for all $0\le t_0<\cdots<t_{n+1}$ and all states $i_0,\dots,i_{n+1}$. Memorylessness is
what buys \cref{eq:markov}: if every pending event is timed by an exponential clock,
the elapsed time carries no information, and the present state is a sufficient summary.

Write $p_{ij}(s,t) = \pr{\xt=j\mid\xs=i}$ for the transition probabilities. Every chain
in this thesis is \emph{time-homogeneous}, meaning that $p_{ij}(s,t)$ depends on $s$ and
$t$ only through the lag $t-s$, so that we may write $p_{ij}(t)$ and collect these into
a matrix $P(t)$. Models are not specified by writing $P(t)$ down, however, but by
writing rates.

\begin{definition}[Infinitesimal generator]
\label{def:generator}
The generator of a time-homogeneous CTMC is the matrix $Q=(q_{ij})$ with
\begin{equation}
\label{eq:generator}
q_{ij} = \lim_{\delt\to0^+}\frac{p_{ij}(\delt)}{\delt}\ \ (i\neq j),
\qquad
q_{ii} = -\sum_{j\neq i}q_{ij},
\end{equation}
whenever these limits exist; equivalently
$p_{ij}(\delt) = \delta_{ij} + q_{ij}\delt + o(\delt)$.
\end{definition}

\noindent
Off-diagonal entries are non-negative and each row sums to zero. The generator also
recovers the pathwise picture of \cref{sec:exp}: while the chain occupies state $i$ it
waits an exponential time of rate $\Lambda_i = -q_{ii} = \sum_{j\neq i}q_{ij}$, and it
then moves to $j$ with probability $q_{ij}/\Lambda_i$. One clock per exit channel, and
\cref{eq:racen} settles the rest. A state $i$ with $q_{ii}=0$ has no exits at all and is
called \emph{absorbing}; in a population model the absorbing state is $i=0$, extinction,
and the whole of this chapter is in one way or another about how a process approaches
it.

Differentiating the Chapman--Kolmogorov relation gives the forward equations,
\begin{equation}
\label{eq:kolmogorov}
\ddt{p_{ij}(t)} = \sum_{k} p_{ik}(t)\,q_{kj},
\qquad\text{that is}\qquad
\dda{P}{t} = P(t)\,Q,
\qquad P(0)=I .
\end{equation}
Conventions for which system is called forward and which backward differ between texts;
\cref{eq:kolmogorov} is the one used throughout, and it is the useful one here because
it reads as a balance of probability flow into and out of state $j$. Written for the
marginal probabilities $p_j(t) = \pr{\xt=j}$ it is the master equation of the model, and
every derivation in \cref{sec:moc} begins by writing it down.

\subsection{Probability generating functions}
\label{sec:pgf}

The master equation is an infinite system of coupled \de s, one per state, and solving
it state by state is rarely possible. The standard remedy is to pack the whole system
into one analytic object.

\begin{definition}[Probability generating function]
\label{def:pgf}
For a \rv{} $\xx$ on $\{0,1,2,\dots\}$ with $\pr{\xx=k}=p_k$, the \pgf{} is
\begin{equation}
\label{eq:pgfdef}
G_{\xx}(z) = \mean{z^{\xx}} = \sum_{k\ge0}p_k z^k,
\qquad |z|\le1 .
\end{equation}
\end{definition}

\noindent
The series converges on the closed unit disc because the coefficients are probabilities,
and the distribution can be read back off it by differentiation. Four consequences are
used constantly:
\begin{equation}
\label{eq:pgfrules}
G_{\xx}(1) = 1,
\qquad
G_{\xx}(0) = p_0,
\qquad
p_k = \frac{1}{k!}G_{\xx}^{(k)}(0),
\qquad
\mean{\xx} = G_{\xx}'(1).
\end{equation}
The second of these is the one that motivates the whole apparatus for small populations:
if $\xx$ is a population size then $G_{\xx}(0)$ is the extinction probability. The
variance follows from the second derivative,
$\text{Var}(\xx) = G_{\xx}''(1)+G_{\xx}'(1)-\bigl(G_{\xx}'(1)\bigr)^2$.

A fifth property does the structural work. If $\xx$ and $\yy$ are independent then
$\mean{z^{\xx+\yy}} = \mean{z^{\xx}}\mean{z^{\yy}}$, so
\begin{equation}
\label{eq:pgfindep}
G_{\xx+\yy}(z) = G_{\xx}(z)\,G_{\yy}(z).
\end{equation}
Independent individuals therefore multiply their generating functions, which is why
generating functions and branching arguments belong together: \cref{eq:pgfindep} is what
turns \enquote{each of the $n$ founders behaves independently} into an $n$th power.

For a process rather than a single \rv{} the generating function acquires a time
argument,
\begin{equation}
\label{eq:pgfprocess}
G(z,t) = \mean{z^{\xt}} = \sum_{j\ge0}p_j(t)\,z^j ,
\end{equation}
and the extinction probability by time $t$ is $p_0(t) = G(0,t)$. Multiplying the master
equation \cref{eq:kolmogorov} by $z^j$, summing over $j$ and recognising the resulting
sums as derivatives of $G$ replaces the infinite system by a single first-order partial
\de{} for $G(z,t)$. That compression is the analytical centre of this chapter. When two
counts must be tracked at once --- particles inside a compartment and particles outside
it, say --- the same construction runs with one variable per count,
\begin{equation}
\label{eq:pgfmulti}
G(x,y,t) = \mean{x^{\xt}y^{\yt}} = \sum_{i\ge0}\sum_{j\ge0}p_{i,j}(t)\,x^iy^j ,
\end{equation}
which is the object solved in \cref{sec:moc} and the natural starting point for the
multi-type models developed later in the thesis.

\subsection{The linear birth--death process}
\label{sec:lbd}

Everything above can now be exercised on the model that all later constructions extend.
Let $\xt$ count individuals, each of which independently gives birth at rate $\lambda$
and dies at rate $\mu$. Because the individuals act independently, the population rates
are linear in the state,
\begin{equation}
\label{eq:lbdrates}
q_{i,i+1} = i\lambda,
\qquad
q_{i,i-1} = i\mu,
\qquad
q_{ii} = -i(\lambda+\mu)
\qquad (i\ge1),
\end{equation}
and $i=0$ is absorbing. The master equation \cref{eq:kolmogorov} reads
\begin{equation}
\label{eq:lbdmaster}
\ddt{p_0} = \mu p_1,
\qquad
\ddt{p_j} = \lambda(j-1)p_{j-1} - (\lambda+\mu)j\,p_j + \mu(j+1)p_{j+1}
\qquad (j\ge1),
\end{equation}
and sequential solution stalls immediately, since each equation involves its
neighbours. Multiplying by $z^j$ and summing, as in \cref{sec:pgf}, collapses the system
to
\begin{equation}
\label{eq:lbdpde}
\pddt{G(z,t)} = (\lambda z-\mu)(z-1)\pdd{G(z,t)}{z},
\qquad
G(z,0) = z^N ,
\end{equation}
for a process started from $\xx_0=N$ individuals. \Cref{eq:lbdpde} is a first-order
linear equation of exactly the type solved by the method of characteristics in
\cref{sec:moc}; here we simply record its solution, which for $\lambda\neq\mu$ is
\begin{equation}
\label{eq:lbdG}
G(z,t) = \lrb{\frac{\mu(1-z) - (\mu-\lambda z)\ee^{(\mu-\lambda)t}}
                   {\lambda(1-z) - (\mu-\lambda z)\ee^{(\mu-\lambda)t}}}^{N}.
\end{equation}
Two checks fix the reading. At $t=0$ the exponentials are unity and the bracket
collapses to $z$, so $G=z^N$; and at $z=1$ numerator and denominator agree, so $G=1$ and
probability is conserved. The mean follows from $\mean{\xt} = \partial_z G|_{z=1}$,
\begin{equation}
\label{eq:lbdmean}
\mean{\xt} = N\ee^{(\lambda-\mu)t},
\end{equation}
which is the deterministic answer and, on its own, an uninformative one: it reports
exponential growth or exponential decay and says nothing about the chance of neither.

That information sits at $z=0$. Setting $z=0$ in \cref{eq:lbdG}, and treating the
critical case $\lambda=\mu$ by taking the limit, gives the probability of extinction by
time $t$ \cite{Allen2010},
\begin{equation}
\label{eq:p0t}
p_0(t) = G(0,t) = \begin{cases}
\lrb{\dfrac{\mu - \mu\ee^{(\mu-\lambda)t}}{\lambda - \mu\ee^{(\mu-\lambda)t}}}^{N},
  & \lambda\neq\mu,\\[12pt]
\lrb{\dfrac{\lambda t}{1+\lambda t}}^{N}, & \lambda=\mu,
\end{cases}
\end{equation}
and the survival probability $S^{(N)}(t) = 1-p_0(t)$. Letting $t\to\infty$ in
\cref{eq:p0t} recovers the classical trichotomy: extinction is certain when
$\lambda\le\mu$, including at $\lambda=\mu$, and occurs with probability
$(\mu/\lambda)^N$ when $\lambda>\mu$. A population with a per-capita growth advantage
can still be lost, and the chance of losing it falls only geometrically in the founding
number $N$.

This process is the continuous-time backbone of the thesis; later chapters add a
catastrophe channel to \cref{eq:lbdrates} and follow the same route from rates to master
equation to generating function. The next section takes the other route, replacing the
exponential clock by a generational one, and asks the same questions of the resulting
discrete-time process.

\FloatBarrier
